\chapter{数据采集与数据处理}

\section{专用数据采集软件}

项目开发了面向 ALOHA 的专用采集软件，而不是临时保存几段视频。当前版本可核验的能力包括：双臂遥操作、连续采集、异步保存、丢弃后重采、键盘/本地脚踏/远程触发、随机起点、回到初始位或不回位两类流程、连续旋转关节、轨迹与视频联合保存、硬件编码不可用时降级、轨迹完整性检查、机器人健康检查以及安全停止。\ev{E02}

\plain{采集软件像一条“数据生产线”：采集员可以连续录制，发现错误就立刻丢弃重来；保存视频的工作在后台完成，避免每录一条都长时间等待。}

本次审计确认的是 2026 年 7 月软件当前能力。由于软件持续演进，不能反推 2026 年 5 月正式模型的所有训练数据都使用过全部新功能。

\section{数据编辑与分析中心}

数据编辑中心集中存放和管理不同机器人数据，不局限于 ALOHA。只读数据库审计得到 51 个有效 ALOHA 项目、2,413 条有效轨迹、2,907,804 帧、58,156.08 秒（约 16.15 小时）；项目汇总数与有效轨迹逐条统计完全一致。另有 562 条轨迹不属于当前有效项目范围，已排除在上述总量外。\ev{E03}

平台已经具备项目与轨迹浏览、单条轨迹与多视角视频查看、标签和批量标签、数据集生成、子任务生成、处理任务跟踪、停止与重新执行等能力，并兼容多类机器人数据。该平台把“散落在机器上的文件”转化为可审查、可版本化和可再次训练的数据资产。

\figfull[.96\textwidth]{data_scope.pdf}{数据资产的三层口径。平台总资产、正式模型唯一训练子集和过滤后可训练数据不能混为同一个数字。}

\section{正式模型实际使用的数据}

\begin{table}[H]
\centering
\caption{正式部署模型训练数据汇总}
\begin{tabularx}{\textwidth}{>{\bfseries}p{34mm}p{38mm}X}
\toprule
指标 & 核验值 & 解释\\
\midrule
唯一数据仓库 & 25 个 & 训练配置含 29 个采样条目，重复条目是采样权重，不是新仓库\\
唯一轨迹 & 1,051 条 & 每条轨迹是一段人工双臂示教\\
唯一帧 & 879,852 帧 & 按声明 50 帧/秒换算约 4.89 小时\\
过滤后可训练帧 & 844,102 帧 & 约 4.69 小时；35,750 帧未进入训练\\
加权采样暴露 & 1,495 条轨迹等效量 & 重点数据被重复抽到，不是新增采集轨迹\\
原始画面 & 4 路，640$\times$480，50帧/秒 & 正式模型只使用顶部、左腕、右腕三路\\
训练画面 & 3 路，224$\times$224 & 训练时统一缩放\\
状态/动作 & 各 14 维 & 双臂各 6 关节和 1 夹爪\\
语言指令 & 有 & 由每个数据仓库的任务描述提供\\
奖励/成功标签 & 无 & 该基础模仿学习数据不能直接计算成功率\\
数据划分 & 仅训练划分 & 没有独立验证集和测试集\\
\bottomrule
\end{tabularx}
\end{table}

公开数据位于 Hugging Face 账户
\href{https://huggingface.co/lyl472324464}{huggingface.co/lyl472324464}。
平台的 Dataset Viewer 支持按行预览、筛选、统计和查看图像/视频字段\cite{hf_dataset_viewer}。本报告用固定版本的元数据、表格和视频文件交叉核验，不只依赖网页标题。

\figfull[.85\textwidth]{sampling_exposure.pdf}{重点场景的重复采样。1,495 是模型在一个采样周期中的轨迹暴露等效量，并非 1,495 条独立采集轨迹。}

\section{数据条件覆盖}

\figfull[.95\textwidth]{condition_coverage.pdf}{正式训练数据中的条件覆盖。类别由数据名称中的明确关键词得到且允许重叠，不能把柱子相加为总量，也不能把数量当作成功率。}

方向变化数据共有 404 条轨迹，无盖条件 158 条，回位条件 135 条，含水条件 136 条，翻转条件 127 条，自由旋转瓶盖条件 111 条。部署训练把自由旋转瓶盖数据重复采样 5 次，说明团队已经主动针对困难条件增加训练暴露。\ev{E04,E11}

\figfull[.92\textwidth]{baseline_episode_length_distribution.pdf}{1,051 条正式训练轨迹的长度分布。中位时长为 12.2 秒，中间 80\% 约为 4.9--28.4 秒；图的极长尾为保证可读性截到 99\% 分位。}

\section{真实训练示教画面}

\figfull[\textwidth]{training_demonstration_keyframes.png}{四类真实训练示教的固定关键帧。每类选择仓库中的中位轨迹，再取 20\%、50\%、80\% 时刻。它们用于说明数据覆盖，不是机器人自主评测结果。}

\section{质量检查、归一化和增强}

对 879,852 条数值记录的全量检查结果如下：

\begin{table}[H]
\centering
\caption{正式训练数据数值质量检查}
\begin{tabularx}{\textwidth}{>{\bfseries}p{42mm}p{28mm}X}
\toprule
检查项 & 结果 & 边界\\
\midrule
状态或动作中的非有限值 & 0 & 全量数值表检查\\
全零状态行 & 0 & 全量数值表检查\\
全零动作行 & 0 & 全量数值表检查\\
轨迹内时间戳倒退/停滞 & 0 条轨迹 & 依据保存时间戳\\
轨迹内帧号不连续 & 0 条轨迹 & 依据保存帧号\\
数值轨迹完全重复组 & 0 组 & 不等于视觉语义无重复\\
损坏视频 & 尚未全量确认 & 只提取并解码了代表性视频\\
训练/测试泄漏 & 无法检查 & 当前只有训练划分\\
\bottomrule
\end{tabularx}
\end{table}

模型对每个状态和动作维度保存第 1 百分位和第 99 百分位。对一个标量 $x$，实际使用的分位数归一化可表示为
\[
\tilde{x}=2\,\frac{\operatorname{clip}(x,q_{01},q_{99})-q_{01}}
{q_{99}-q_{01}}-1,
\]
即把大多数正常值映射到 $[-1,1]$，同时限制极端值的影响；输出动作再按相反方向还原。\ev{E09} 训练中使用图像随机裁剪、缩放、旋转与颜色扰动，但是否对所有数据仓库以完全相同概率生效，应以当次训练记录为准。

\limitbox{当前数据只有训练划分，没有独立验证和测试集；视频未逐帧全部解码；瓶盖存在性、瓶身朝向和任务阶段也没有逐帧人工标签。这三项是下一年度建立可信评测的直接阻塞。}

\begin{table}[H]
\centering
\caption{数据审计结论：可以确认与不能确认}
\begin{tabularx}{\textwidth}{>{\bfseries}p{33mm}X}
\toprule
证据等级 & 结论\\
\midrule
可以确认 & 数据规模、固定版本、轨迹长度、三路训练相机、14维状态/动作、数值完整性、重点采样和任务指令均可追溯。\\
只能部分确认 & 多条件名称说明采集意图，但没有逐帧语义标注；代表视频能解码，但未对全部视频逐帧检查。\\
不能确认 & 没有独立验证/测试，无法证明没有数据泄漏；没有成功标签，无法从训练数据直接计算自主任务成功率。\\
\bottomrule
\end{tabularx}
\end{table}
